% =====================================================================
% cap-03-motor-matematico.tex — Capítulo 4: Motor Matemático
% =====================================================================

\chapter{Motor Matemático}
\label{cap:motor}

Este capítulo formaliza o \textit{pipeline} matemático que transforma
uma nuvem bruta de pontos do sensor Kinect em uma malha de quadrados
coloridos sobre a areia. Cada uma das quatro etapas é apresentada com
rigor algébrico e relacionada à sua implementação no módulo
\texttt{motor\_caixao\_areia.py}.

\section{Passo 1 — Ajuste de Plano: RANSAC seguido de SVD}
\label{sec:passo1}

Dada uma nuvem de $N$ pontos $\{\mathbf{p}_i\}_{i=1}^{N} \subset
\mathbb{R}^{3}$ capturados pelo Kinect sobre a superfície plana de
referência (a tampa apoiada sobre o caixão, ou a base de madeira vazia),
o problema é determinar o plano $\Pi: ax + by + cz + d = 0$ que melhor
se ajusta a essa superfície. O campo de visão do Kinect, contudo, é mais
largo que a área útil do caixão: a nuvem capturada inclui também a
moldura de madeira, o piso ao redor e ruído da sala — pontos que não
pertencem ao plano de referência e que distorceriam um ajuste de
mínimos quadrados ingênuo. Por isso a calibração oficial
(\texttt{main.\_executar\_calibracao}) combina dois métodos em
sequência: primeiro o RANSAC isola o maior subconjunto de
pontos coplanares (Seção~\ref{subsec:ransac}); em seguida, o
SVD refina a normal com precisão sub-milimétrica, mas apenas
sobre esse subconjunto de \emph{inliers}.

\subsection{Isolamento de inliers por RANSAC}

A função \texttt{ajustar\_plano\_ransac()} executa $n_{\text{iter}} =
1000$ iterações por padrão. A cada iteração, sorteia-se uma trinca de
pontos, constrói-se o plano candidato que passa por eles e contam-se os
\emph{inliers}: pontos da nuvem completa cuja distância ao plano
candidato é inferior a $d_{\text{in}} = 0{,}03$~m (3~cm). O plano
vencedor é o que acumula o maior número de inliers ao longo das
iterações. Se a fração de inliers do plano vencedor for inferior a
$30\,\%$ da nuvem, a função levanta \texttt{RuntimeError} — sintoma de
que a superfície de referência não domina o campo de visão do sensor
(tampa mal posicionada, por exemplo).

\subsection{Refinamento por SVD}

A solução clássica de mínimos quadrados procede em três etapas sobre o
conjunto de inliers identificado pelo RANSAC (e apenas sobre ele).
Primeiro, calcula-se o centroide e a matriz centralizada
\begin{equation}
\bar{\mathbf{p}} = \frac{1}{N} \sum_{i=1}^{N} \mathbf{p}_i,
\qquad
M = \begin{bmatrix}
(\mathbf{p}_1 - \bar{\mathbf{p}})^{T} \\
\vdots \\
(\mathbf{p}_N - \bar{\mathbf{p}})^{T}
\end{bmatrix} \in \mathbb{R}^{N \times 3}.
\end{equation}

Em seguida, aplica-se a decomposição em valores singulares
$M = U \Sigma V^{T}$ com $\Sigma = \text{diag}(\sigma_1, \sigma_2,
\sigma_3)$ e $\sigma_1 \geq \sigma_2 \geq \sigma_3 \geq 0$. O vetor
normal procurado é a última linha de $V^{T}$:
\begin{equation}
\mathbf{n} = V^{T}_{[3,:]}.
\end{equation}

Finalmente, o coeficiente $d$ é calculado pela substituição do
centroide na equação do plano: $d = -\mathbf{n} \cdot \bar{\mathbf{p}}$.
Por convenção, exige-se que $n_z > 0$ (vetor apontando para cima, em
direção ao Kinect); caso contrário, $\mathbf{n}$ é negado.

A função \texttt{ajustar\_plano\_svd()} encapsula esse procedimento e é
coberta por quatro testes unitários que verificam (i) unicidade da
normal, (ii) plano horizontal $z = 5$, (iii) plano inclinado e
(iv) exceção para $N < 3$. A função \texttt{ajustar\_plano\_ransac()},
que aplica o RANSAC descrito acima e delega o refinamento a
\texttt{ajustar\_plano\_svd()} sobre os inliers, é coberta por cinco
testes unitários adicionais (classe \texttt{TestRANSAC}), incluindo o
cenário motivador da especificação — uma nuvem sintética com moldura e
piso como \textit{outliers} ao redor da tampa — e a verificação de que
o refinamento final usa exclusivamente os inliers selecionados.

\section{Passo 2 — Construção da Base da Mesa}
\label{sec:passo2}

A normal $\mathbf{n}$ define inequivocamente o eixo $Z_{\text{mesa}}$,
mas a base ortonormal completa requer dois eixos adicionais $X_{\text{mesa}}$
e $Y_{\text{mesa}}$ ortogonais entre si e a $\mathbf{n}$. A escolha
desses eixos é livre a menos de rotação em torno da normal; adotamos a
convenção mais simples possível, baseada em uma semente padrão.

O algoritmo é o seguinte:

\begin{enumerate}
    \item $Z_{\text{mesa}} \leftarrow \mathbf{n} / \|\mathbf{n}\|$;
    \item Escolhe-se uma semente $\mathbf{s} = (1, 0, 0)^{T}$. Se
          $|\mathbf{s} \cdot Z_{\text{mesa}}| \geq 0{,}9$ (semente
          quase paralela à normal), substitui-se por $(0, 1, 0)^{T}$;
    \item Aplica-se Gram-Schmidt:
          $X_{\text{mesa}} = \dfrac{\mathbf{s} - (\mathbf{s} \cdot Z_{\text{mesa}}) Z_{\text{mesa}}}{\| \mathbf{s} - (\mathbf{s} \cdot Z_{\text{mesa}}) Z_{\text{mesa}} \|}$;
    \item Calcula-se $Y_{\text{mesa}} = Z_{\text{mesa}} \times
          X_{\text{mesa}}$, garantindo automaticamente um sistema
          dextrogiro.
\end{enumerate}

A ortonormalidade do conjunto $\{X_{\text{mesa}}, Y_{\text{mesa}},
Z_{\text{mesa}}\}$ — isto é, $X \cdot Y = X \cdot Z = Y \cdot Z = 0$ e
$\|X\| = \|Y\| = \|Z\| = 1$ — é assegurada por testes unitários
dedicados (\texttt{TestConstruirBase}).

\section{Passo 3 — Transformação Afim 4$\times$4}
\label{sec:passo3}

A transformação $T \in \mathbb{R}^{4 \times 4}$ que leva pontos do
referencial do Kinect ao referencial da mesa empilha os eixos da mesa
como linhas da matriz de rotação:
\begin{equation}
R = \begin{bmatrix}
X_{\text{mesa}}^{T} \\
Y_{\text{mesa}}^{T} \\
Z_{\text{mesa}}^{T}
\end{bmatrix},
\qquad
\mathbf{t} = -R \, \bar{\mathbf{p}},
\qquad
T = \begin{bmatrix} R & \mathbf{t} \\ \mathbf{0}^{T} & 1 \end{bmatrix}.
\end{equation}

\noindent A propriedade fundamental dessa construção é que, para
qualquer ponto $\mathbf{p}$ pertencente ao plano da mesa, a coordenada
transformada $z_{\text{mesa}} = 0$ — isto é, o plano da mesa torna-se o
plano $XY$ do novo sistema. Essa invariante é verificada pelo teste
\texttt{TestMatrizTransformacao}.

\subsection{Composição com $T_{\text{shift}}$ — centro lógico e cota zero na base}
\label{subsec:T-shift}

A matriz $T$ acima leva o centroide $\bar{\mathbf{p}}$ do plano
detectado à origem $(0, 0, 0)$. Entretanto, a convenção interna
de discretização e renderização (Capítulo~\ref{cap:rendering}) adota a
mesa como o domínio $[0, L_x] \times [0, L_y] \times [0, h_{\max}]$, com
a origem $XY$ em um canto e a cota $Z = 0$ fixada na base
de madeira do caixão (alturas de areia sempre positivas, crescendo
para cima). Sem ajuste em $XY$, metade dos pontos da nuvem (aqueles com
$x < 0$ ou $y < 0$) seria projetada fora do domínio e colapsada
pelo \texttt{clip} nas células de borda da grade — um defeito
observado em ensaio preliminar de campo, no qual a projeção exibia
apenas um pequeno retângulo correspondente ao quadrante positivo. E sem
ajuste em $Z$, a cota zero coincidiria com o plano de calibração
em vez da base do caixão — o que é apenas correto quando esse plano é a
própria base (modo base vazia), mas não quando ele é a tampa
apoiada $h_{\max}$ acima dela (modo tampa, oficial).

A solução, identificada e corrigida durante a fase de validação, é
compor $T$ com uma translação $T_{\text{shift}}$ que (i) desloca o
centroide para o centro lógico da mesa em $XY$ e (ii) desce a origem em
$Z$ até a base do caixão sempre que o plano de calibração for a tampa:
\begin{equation}
\label{eq:T-shift}
T_{\text{shift}} =
\begin{bmatrix}
1 & 0 & 0 & L_x / 2 \\
0 & 1 & 0 & L_y / 2 \\
0 & 0 & 1 & \delta_{\text{tampa}} \, h_{\max} \\
0 & 0 & 0 & 1
\end{bmatrix},
\qquad
T_{\text{final}} = T_{\text{shift}} \cdot T,
\end{equation}

\noindent onde $\delta_{\text{tampa}} = 1$ no modo tampa e
$\delta_{\text{tampa}} = 0$ no modo base vazia. Após essa composição, o
ponto físico que se encontra exatamente sob o sensor (centroide do
plano detectado) cai em $(L_x/2,\, L_y/2,\, \delta_{\text{tampa}}
h_{\max})$: no modo tampa, na cota $h_{\max}$ (o topo do caixão); no
modo base, na cota $0$ (o fundo, que é o próprio plano de
calibração). Em ambos os casos, a nuvem ocupa naturalmente $[0, L_x]
\times [0, L_y]$ em $XY$ e a base de madeira corresponde sempre a
$Z_{\text{mesa}} = 0$, independentemente do modo de calibração
escolhido. Adicionalmente, a função \texttt{discretizar\_nuvem\_em\_grade()}
aplica um filtro \texttt{dentro} que descarta pontos cujas
coordenadas $XY$ extrapolem o domínio físico da mesa, o que é típico do
campo de visão do Kinect montado a 2,5~m (FOV horizontal de
aproximadamente 70° cobre uma área maior do que a mesa).

Como validação de sanidade, antes de compor $T_{\text{shift}}$
o sistema compara a distância medida sensor–plano ($-\bar{p}_z$, uma
vez que $Z$ cresce em direção ao sensor no referencial original) com a
distância informada pelo operador na interface gráfica — a distância ao
plano da tampa no modo tampa, ou essa mesma distância somada a
$h_{\max}$ no modo base. Divergências acima de 15~cm abortam a
calibração com uma mensagem diagnóstica, protegendo contra tampa mal
posicionada ou erro de digitação nos campos físicos da mesa.

\section{Passo 4 — Projeção 3D$\to$2D pelo Modelo de Tsai}
\label{sec:passo4}

Os vértices da malha de quadrados (Capítulo~\ref{cap:rendering}) são
projetados do referencial 3D da mesa para o plano 2D do projetor pela
fórmula geral~\eqref{eq:pinhole}, com distorções óticas residuais
$(k_1, k_2, p_1, p_2, k_3)$ encapsuladas pela função
\texttt{cv2.projectPoints} do OpenCV.

Em modo simulação e em modo real, os parâmetros
intrínsecos e extrínsecos são configurados de forma a mapear o domínio
físico $[0, L_x] \times [0, L_y]$ na janela de projeção. Adotando uma
distância virtual da câmera $d_{\text{cam}} = 10$~m, $\mathbf{t}_{\text{ext}} =
(0, 0, d_{\text{cam}})^{T}$ e $R_{\text{ext}} = I_3$, obtém-se
\begin{equation}
\label{eq:fxfy-virtual}
f_x = \frac{W \cdot d_{\text{cam}}}{L_x}, \qquad
f_y = \frac{H \cdot d_{\text{cam}}}{L_y},
\qquad (c_x, c_y) = (0, 0),
\end{equation}

\noindent onde $(W, H)$ é a resolução do projetor. Substituindo em
\eqref{eq:pinhole}, vê-se que o ponto $(x, y, 0)$ projeta-se exatamente
em
\begin{equation}
u = \frac{W}{L_x} \, x, \qquad v = \frac{H}{L_y} \, y,
\end{equation}

\noindent o que é precisamente o mapeamento desejado de
$[0, L_x] \times [0, L_y]$ para $[0, W] \times [0, H]$. Essa parametrização
é compartilhada pelos dois modos de operação, o que simplifica a
calibração e elimina assimetrias entre simulação e hardware real. Como
o projetor é montado fisicamente alinhado — apontando reto para baixo,
centrado sobre a mesa —, esses parâmetros podem ser deduzidos
analiticamente, sem necessidade de um procedimento de estimação
experimental.

A Tabela~\ref{tab:intrinsecos-projetor} instancia a
Equação~\eqref{eq:fxfy-virtual} com os valores efetivamente usados
pelo sistema: a resolução de projeção padrão
(\texttt{RESOLUCAO\_PROJETOR}, fixa no código) e a distância virtual
$d_{\text{cam}} = 10$~m adotada em \texttt{main.py}, aplicadas à mesa
física de referência deste trabalho ($1,5~\text{m} \times
1,5~\text{m}$, Capítulo~\ref{cap:rendering}).

\begin{table}[!ht]
\centering
\caption{Parâmetros intrínsecos da câmera virtual do projetor (modelo analítico da Equação~\eqref{eq:fxfy-virtual}), instanciados para a mesa física de referência deste trabalho.}
\label{tab:intrinsecos-projetor}
\begin{tabular}{ll}
\toprule
\textbf{Parâmetro} & \textbf{Valor} \\
\midrule
Resolução do projetor $(W, H)$          & $640 \times 480$ pixels \\
Distância virtual da câmera $d_{\text{cam}}$ & $10{,}00$ m \\
Dimensões da mesa $(L_x, L_y)$          & $1{,}50 \times 1{,}50$ m \\
Distância focal $f_x$                   & $4\,266{,}67$ pixels \\
Distância focal $f_y$                   & $3\,200{,}00$ pixels \\
Centro óptico $(c_x, c_y)$              & $(0{,}00,\; 0{,}00)$ pixels \\
\bottomrule
\end{tabular}
\end{table}

Ao contrário dos parâmetros intrínsecos do Kinect
(Tabela~\ref{tab:intrinsecos}), que são uma especificação fixa do
fabricante do sensor, $f_x$ e $f_y$ do projetor não são uma constante
de hardware: dependem das dimensões físicas da mesa configuradas pelo
operador na inicialização do sistema (Apêndice~\ref{ap:roteiro}), já
que a Equação~\eqref{eq:fxfy-virtual} os deriva justamente dessas
dimensões. Os valores acima instanciam essa equação para a mesa de
referência ($1,5~\text{m} \times 1,5~\text{m}$) usada em todos os
ensaios e testes deste trabalho; $W$, $H$ e $d_{\text{cam}}$, por sua
vez, são constantes do sistema, fixas independentemente da mesa
configurada.

\subsection{Calibração empírica alternativa do projetor (tabuleiro de xadrez)}
\label{subsec:calib-empirica-projetor}

Caso a premissa de alinhamento físico não se sustente — projetor
montado com alguma inclinação residual, por exemplo —, o sistema
disponibiliza um caminho alternativo de estimação empírica dos
parâmetros de Tsai, no qual a etapa de otimização que o artigo original
de Tsai~\citeonline{tsai1987versatile} descreve é de fato executada. A
função \texttt{gerar\_imagem\_tabuleiro()} produz um padrão de
tabuleiro de xadrez com cantos internos em coordenadas conhecidas do
projetor; essa imagem é projetada temporariamente na janela
\texttt{Projecao\_Areia} sobre a superfície de referência ainda no
lugar (Seção~\ref{sec:saidas}). Uma câmera RGB captura o tabuleiro
projetado; \texttt{encontrar\_cantos\_tabuleiro()} localiza os cantos
na imagem capturada via \texttt{cv2.findChessboardCorners}; e
\texttt{cantos\_rgb\_para\_pontos\_mesa()} converte cada canto para uma
coordenada 3D no referencial da mesa, por \textit{back-projection}
(Equação~\eqref{eq:backproj}) seguida de $T_{\text{final}}$. Com as
correspondências 3D~$\leftrightarrow$~2D resultantes,
\texttt{calibrar\_projetor()} invoca \texttt{cv2.calibrateCamera} para
estimar $f_x, f_y, c_x, c_y$ e os coeficientes de distorção por
otimização não linear, a partir de um chute inicial de focal
configurável (fator ``\textit{throw ratio}'', padrão $1{,}2 \times$ a
largura do projetor em pixels — exposto ao operador por um controle
deslizante na GUI e por um \textit{trackbar} na janela
\texttt{Gabarito\_MDE} durante a calibração, permitindo corrigi-lo caso
a otimização convirja para um mínimo local incorreto). Esse caminho
alternativo é exercitado pelas classes de teste
\texttt{TestCalibracaoProjetorOpenCV}, \texttt{TestGeracaoTabuleiro},
\texttt{TestCantosParaPontosMesa} e
\texttt{TestPipelineCalibracaoProjetor} (Capítulo~\ref{cap:resultados}),
mas permanece inativo por padrão: a configuração corrente do
sistema opera com os parâmetros analíticos da
Equação~\eqref{eq:fxfy-virtual}, reservando a rotina empírica para o
cenário de montagem não idealizada.

\section{Discussão — RANSAC e SVD: uma composição deliberada}
\label{sec:discussao-svd}

Uma versão anterior deste trabalho empregava SVD puro, sem RANSAC, para
o ajuste de plano, sob a premissa de que a nuvem de calibração seria
capturada já livre de \textit{outliers}. A experimentação em
campo mostrou que essa premissa não se sustenta: o campo de visão do
Kinect~v2 é mais largo que a área útil do caixão de areia, de modo que
qualquer captura da superfície de referência inclui também a moldura de
madeira do caixão, o piso ao redor e ruído da sala. Um SVD aplicado
ingenuamente sobre a nuvem inteira minimiza a soma dos quadrados das
distâncias de \emph{todos} os pontos, inclusive desses outliers, o que
distorce sistematicamente a normal calculada. A solução adotada — e
validada nos ensaios de campo relatados no
Capítulo~\ref{cap:resultados} — não é substituir o SVD pelo RANSAC, mas
compor os dois, cada um explorando a etapa em que é mais forte:

\begin{itemize}
    \item O RANSAC (Seção~\ref{subsec:ransac}) é robusto a
          \textit{outliers} por construção: ele isola o maior
          subconjunto de pontos coplanares antes de qualquer ajuste
          numérico, funcionando mesmo com até $70\,\%$ da nuvem
          ocupada por pontos que não pertencem ao plano de referência
          (limite imposto por \texttt{min\_inliers\_ratio}~$= 0{,}3$);
    \item O SVD, aplicado apenas sobre os inliers
          selecionados pelo RANSAC, produz o estimador de mínimos
          quadrados \emph{ótimo} para esse subconjunto — mais preciso
          do que o plano obtido a partir de apenas três pontos
          amostrados em uma única iteração do RANSAC;
    \item A solução SVD é \emph{única} (a menos do sinal da normal,
          tratado por convenção), o que facilita a verificação por
          testes unitários com valores esperados \emph{hardcoded}
          sobre nuvens sintéticas 100\,\% coplanares — cenário em que
          o parâmetro \texttt{usar\_ransac=False} permanece disponível
          e suficiente, dispensando o custo das iterações RANSAC;
    \item Em contrapartida, o SVD isolado seria \emph{insuficiente}
          para o cenário real da tampa: sem a etapa de RANSAC, a
          normal calculada é sistematicamente puxada na direção dos
          outliers, produzindo o mesmo artefato de projeção deslocada
          relatado nos ensaios preliminares de campo
          (Seção~\ref{subsec:T-shift}).
\end{itemize}

\noindent A escolha entre os dois regimes é, portanto, guiada pela
origem dos dados: SVD puro para nuvens sintéticas já garantidamente
coplanares (a maioria dos testes unitários do motor matemático); RANSAC
seguido de SVD para qualquer nuvem capturada por hardware real, na qual
a presença de outliers geométricos é a regra, não a exceção.

\noindent Essa decisão é validada empiricamente: a suíte de 83 testes
unitários do motor matemático passa integralmente, incluindo os cinco
casos dedicados especificamente ao RANSAC (\texttt{TestRANSAC}) e os
oito casos da calibração pela tampa (\texttt{TestCalibracaoTampa}), com
as constantes algébricas explicitadas na própria suíte.
